% from ega4-errata-addenda-list3-fr.tex
\oldpage[IV]{360}
$f(T_\alpha)=Y$ et $f|T_\alpha$ est quasi-compact. On peut par suite
appliquer à $f|T_\alpha$ le résultat de I), et il existe un diviseur
$D_\alpha$ sur $Y$ tel que
\[
 (f|T_\alpha)^*(D_\alpha)=D'|T_\alpha;
\]
en vertu de la transitivité des images réciproques de diviseurs par
morphismes plats {\hyperlink{ega4.c04.statement.21-4-4}{(21.4.4)}}, on
a aussi $(f|T)^*(D_\alpha)=D'|T$. On en conclut que pour deux indices
$\alpha$, $\beta$ distincts, on a
\[
 (f|T)^*(D_\alpha)=(f|T)^*(D_\beta),
\]
d'où (par {\hyperlink{ega4.c04.prop.21-4-8}{(21.4.8)}}),
$D_\alpha=D_\beta$. Si $D$ est la valeur commune des $D_\alpha$, on a
donc $f^*(D)|T_\alpha=D'|T_\alpha$ pour tout $\alpha$, et par suite
$f^*(D)=D'$.
\end{proof}

\begin{corollary}[21.4.10]
\EGAReaderTarget{ega4.c04.cor.21-4-10}{IV.21.4.10}
Supposons vérifiées les hypothèses générales de
{\hyperlink{ega4.c04.prop.21-4-9}{(21.4.9)}}, et en outre supposons que
$Y$ soit normal et que pour tout $y\in Y$ tel que
$\dim(\sh{O}_{Y,y})=1$, la fibre $X_y=f^{-1}(y)$ n'ait pas de cycle
premier associé immergé. Pour qu'un diviseur $D'$ sur $X$ soit de la
forme $f^*(D)$, où $D\in\operatorname{Div}(Y)$, il faut et il suffit
que les deux conditions suivantes soient vérifiées :
\begin{enumerate}[label=\textup{\arabic*\textdegree}]
\item pour tout point maximal $y$ de $Y$, l'image réciproque de $D'$
dans $\operatorname{Div}(X_y)$ soit nulle;
\item pour tout $y\in Y$ tel que $\dim(\sh{O}_{Y,y})=1$, il existe un
entier $n_y$ tel que, pour tout point maximal $x$ de $X_y$, l'image
réciproque de $D'$ dans
$\operatorname{Div}(\operatorname{Spec}(\sh{O}_{X,x}))$ soit égale à
$\operatorname{div}(t_y^{n_y})$, où $t_y$ est l'image par
l'homomorphisme $\sh{O}_{Y,y}\to\sh{O}_{X,x}$ d'une uniformisante de
l'anneau de valuation discrète $\sh{O}_{Y,y}$.
\end{enumerate}
\end{corollary}

\begin{proof}
Appliquons le critère {\hyperlink{ega4.c04.prop.21-4-9}{(21.4.9)}}.
Aux points maximaux $y$ de $Y$, on a $\operatorname{Div}(Y_1)=0$
puisque $\sh{O}_{Y,y}$ est artinien, et la condition 1\textdegree{} de
(21.4.10) est équivalente au critère de (21.4.9) en un tel point. Il
faut voir que les conditions 1\textdegree{} et 2\textdegree{} impliquent
de même la condition de (21.4.9) en un point $y\in Y$ où
$\operatorname{prof}(\sh{O}_{Y,y})=1$; comme $Y$ est normal, ces points
sont aussi ceux où $\dim(\sh{O}_{Y,y})=1$, ou encore où
$\sh{O}_{Y,y}$ est un anneau de valuation discrète
{\hyperref[canon-index.ega4.c04.thm.5-8-6]{(5.8.6)}}. On peut donc se borner au
cas où $Y=\operatorname{Spec}(\sh{O}_{Y,y})$. Si $t$ est une
uniformisante de $\sh{O}_{Y,y}$, l'application
$n\mapsto\operatorname{div}(t^n)$ est un isomorphisme de $\mathbb Z$
sur $\operatorname{Div}(Y)$; la condition de (21.4.9) signifie que
$D'$ est de la forme $\operatorname{div}(t^n\circ f)$ pour un
$n\in\mathbb Z$. Cela montre aussitôt que cette condition entraîne la
condition 2\textdegree{} de l'énoncé. Inversement, supposons cette
dernière vérifiée (ainsi que la condition 1\textdegree{}); posant pour
simplifier $n=n_y$, tout revient à voir que l'on a alors
$D'=\operatorname{div}(t^n\circ f)$. Si $\eta$ est le point générique
de $Y$, les germes de $D'$ et de $\operatorname{div}(t^n\circ f)$ en
un point quelconque de $X_\eta$ sont tous deux nuls en vertu de
1\textdegree{}. De même, les germes de $D'$ et
$\operatorname{div}(t^n\circ f)$ en un point maximal de $X_y$ sont
égaux par la condition 2\textdegree{}. Mais en un point $x\in X_y$ qui
n'est pas maximal, on a
$\operatorname{prof}(\sh{O}_{X_y,x})\geq1$ puisque $X_y$ n'a pas de
cycle premier associé immergé, donc aussi
$\operatorname{prof}(\sh{O}_{X,x})\geq2$ par
{\hyperref[canon-index.ega4.c04.prop.6-3-1]{(6.3.1)}}. On a donc
$D'=\operatorname{div}(t^n\circ f)$ en vertu de
{\hyperlink{ega4.c04.prop.21-1-8}{(21.1.8)}}.
\end{proof}

\begin{corollary}[21.4.11]
\EGAReaderTarget{ega4.c04.cor.21-4-11}{IV.21.4.11}
Supposons vérifiées les hypothèses générales de
{\hyperlink{ega4.c04.prop.21-4-9}{(21.4.9)}}, et supposons en outre que
$Y$ soit normal et que pour tout $y\in Y$ tel que
$\dim(\sh{O}_{Y,y})=1$, la fibre $X_y$ soit intègre. Alors, pour qu'un
diviseur $D'$ soit de la forme $f^*(D)$, où
$D\in\operatorname{Div}(Y)$, il faut et il suffit que, pour tout point
maximal $y$ de $Y$, l'image réciproque de $D'$ dans
$\operatorname{Div}(X_y)$ soit nulle.
\end{corollary}

\begin{proof}
Il suffit de vérifier que la condition 2\textdegree{} de
{\hyperlink{ega4.c04.cor.21-4-10}{(21.4.10)}} est alors
automatiquement vérifiée. Or, si $\dim(\sh{O}_{Y,y})=1$ et si $X_y$ est
intègre, on a, en l'unique point maximal $x$ de $X_y$,
$\dim(\sh{O}_{X,x})=1$ en vertu de
{\hyperref[canon-index.ega4.c04.cor.6-1-3]{(6.1.3)}}; en outre, comme
$\sh{O}_{Y,y}$ est un anneau régulier et $\sh{O}_{X_y,x}$ un corps,
$\sh{O}_{X,x}$ est régulier
{\hyperref[canon-index.ega4.c04.prop.6-5-1]{(6.5.1)}}, donc un anneau de valuation
discrète, et si $t$ est une uniformisante de $\sh{O}_{Y,y}$, son image
$t_y$ dans $\sh{O}_{X,x}$ est une uniformisante de $\sh{O}_{X,x}$
($0$, 17.3.3); l'image réciproque de $D'$ dans
$\operatorname{Div}(\operatorname{Spec}(\sh{O}_{X,x}))$ est donc bien
de la forme

